\chapter{Background and Related Work}
\label{chapter:background}

This chapter outlines the framework motivating this internship by contrasting classical policy gradients with the challenges of mean-field control. We review finite-horizon Markov decision processes and compare model-based, value-based and policy-based approaches. We then introduce likelihood-ratio policy gradients and REINFORCE, and highlight a defining structural assumption: environment transitions are completely independent of policy parameters given a state--action pair.

Mean-field control breaks this assumption. Because the environment relies on a population distribution driven by the policy, it creates a feedback loop that directly invalidates the classical likelihood-ratio trick. By identifying exactly where this breakdown occurs, we motivate the introduction of a perturbation framework designed to recover a model-free score representation.

\section{Markov decision processes}
\label{sec:background-mdp}

Markov decision processes (MDPs) provide the standard mathematical framework for sequential decision-making under uncertainty and form the basis of classical reinforcement learning \cite{puterman1994mdp, sutton2018reinforcement}. An agent interacts with a stochastic environment over a sequence of time steps: at each time, it observes the current state, selects an action according to a policy, receives a reward, and the system evolves randomly to a new state.

Throughout this report, we consider a finite time horizon $\Tc = \{0,\ldots,T\}$, although the same concepts could be extended to infinite-horizon and discounted settings. Let $\Xc$ denote the state space and $\Ac$ the action space. These spaces may be finite, as in the discrete-state experiments considered later, or more general measurable spaces, as required for continuous-state problems. The model is specified by an initial distribution $\mu_0\in\Pc(\Xc)$, transition kernels $P_t(\d x'\mid x,a)$, running costs $r_t:\Xc\times\Ac\to\R$, and a terminal cost $g:\Xc\to\R$.

A randomized Markov policy is a sequence $\pi=(\pi_t)_{t=0}^{T-1}$ of kernels from $\Xc$ to $\Ac$. Under $\pi$, the controlled process satisfies
\begin{align}
    \label{eq:background-mdp_dynamics}
    \begin{cases}
        X_0 &\sim \mu_0,\\
        \alpha_t &\sim \pi_t(\cdot\mid X_t),\\
        X_{t+1} &\sim P_t(\cdot\mid X_t,\alpha_t), \qquad t=0,\ldots,T-1.
    \end{cases}
\end{align}
A deterministic feedback rule is recovered when $\pi_t(\cdot\mid x)$ is a Dirac measure. Randomized policies are used here because their action probabilities can be differentiated, which is the basis of the likelihood-ratio estimator introduced below. They also provide the exploration needed to compare different actions from simulated trajectories.

The last relation encodes the Markov property: conditionally on $(X_t,\alpha_t)$, the law of $X_{t+1}$ is independent of the preceding trajectory. Consequently, the law of $\tau=(X_0,\alpha_0,\ldots,X_{T-1},\alpha_{T-1},X_T)$ factorizes as
\begin{align}
    \label{eq:background-mdp_trajectory_law}
    \P^\pi(\d\tau) = \mu_0(\d x_0) \prod_{t=0}^{T-1}\pi_t(\d a_t\mid x_t) P_t(\d x_{t+1}\mid x_t,a_t).
\end{align}

The performance criterion is the expected cumulative cost
\begin{align}
    \label{eq:background-mdp_objective}
    J(\pi) = \E^\pi\left[\sum_{t=0}^{T-1} r_t(X_t,\alpha_t) + g(X_T) \right], \qquad \inf_\pi J(\pi).
\end{align}
When the model is known, this problem can be approached by dynamic programming. Reinforcement learning instead assumes access to sampled transitions and costs, without requiring the kernels to be known or tractable in closed form. For this work, the important consequence of \eqref{eq:background-mdp_trajectory_law} is that the environment kernels do not change when the parameters of a policy are varied. This separation makes the classical likelihood-ratio policy gradient model-free.

Reinforcement learning methods differ in what they attempt to learn from these samples. Value-based methods first estimate a value or state-action value function and then derive a policy. Policy-gradient methods instead parameterize the policy and optimize its performance directly. The latter choice is natural for this work because the policy is already the object being parametrized, whereas the transition mechanism is treated as a simulator. It also avoids differentiating the dynamics, at the cost of a noisier Monte Carlo gradient.

\section{Policy-gradient methods}
\label{sec:background-policy_gradient}

Policy-gradient methods optimize a differentiable family of randomized policies directly. They include the REINFORCE algorithm of \cite{williams1992reinforce} and the broader policy-gradient framework of \cite{sutton1999policygradient}. Let $\Theta\subseteq \R^D$ and suppose that $\pi_t^\theta$ admits a positive differentiable density with respect to a fixed reference measure $\nu_\Ac$: $\pi_t^\theta(\d a\mid x) = p_t^\theta(a\mid x)\nu_\Ac(\d a)$.
We write $J(\theta)=J(\pi^\theta)$ and seek a trajectory-based estimator of $\nabla_\theta J(\theta)$.

\subsection{Likelihood-ratio representation}

The score-function, or likelihood-ratio, identity states that, under conditions justifying differentiation under the integral sign,
\begin{align}
    \label{eq:background-likelihood_ratio_identity}
    \nabla_\theta \E_{Z\sim q_\theta}[F(Z)] = \E_{Z\sim q_\theta}\left[F(Z)\nabla_\theta\log q_\theta(Z) \right]
\end{align}
whenever $F$ has no explicit dependence on $\theta$ \cite{lecuyer1990likelihood}. Applied to the trajectory law \eqref{eq:background-mdp_trajectory_law}, the terms involving $\mu_0$ and $P_t$ disappear from the score because they do not depend on $\theta$. Hence
\begin{align}
    \label{eq:background-trajectory_score}
    \nabla_\theta\log q_\theta(\tau) = \sum_{t=0}^{T-1} \nabla_\theta\log p_t^\theta(a_t\mid x_t).
\end{align}
The transition kernels still determine which trajectories are observed, but their values and derivatives are not needed in the estimator.

Let $C_t = \sum_{s=t}^{T-1} r_s(X_s,\alpha_s) + g(X_T)$ denote the cumulative cost-to-go from time $t$. The policy score has zero conditional mean,
\begin{align}
    \label{eq:background-score_zero_expectation}
    \E\left[\nabla_\theta\log p_t^\theta(\alpha_t\mid X_t) \mid X_t \right] = \int_\Ac \nabla_\theta p_t^\theta(a\mid X_t)\nu_\Ac(\d a) = 0.
\end{align}
Costs incurred before time $t$ are measurable before $\alpha_t$ is sampled and therefore have zero expectation when multiplied by the score. Combining this causality argument with \eqref{eq:background-likelihood_ratio_identity} gives the finite-horizon policy-gradient formula
\begin{align}
    \label{eq:background-reinforce_gradient}
    \nabla_\theta J(\theta) = \E^{\pi^\theta} \left[\sum_{t=0}^{T-1} \nabla_\theta\log p_t^\theta(\alpha_t\mid X_t) C_t \right].
\end{align}
Thus credit is assigned using future costs only, rather than the full
trajectory cost.

\subsection{Monte Carlo estimator and baselines}

Given $N$ independent trajectories generated under $\pi^\theta$, let
$C_t^{(k)}$ be the cost-to-go along trajectory $k$. Equation \eqref{eq:background-reinforce_gradient} gives the REINFORCE estimator
\begin{align}
    \label{eq:background-reinforce_estimator}
    \widehat{\nabla_\theta J}(\theta) = \frac1N \sum_{k=1}^N \sum_{t=0}^{T-1} \nabla_\theta\log p_t^\theta(\alpha_t^{(k)}\mid X_t^{(k)}) C_t^{(k)}.
\end{align}
In this sense, REINFORCE replaces the unknown conditional expected cost-to-go by its realized Monte Carlo value on each trajectory. This avoids learning a separate value function, but is also the main source of variance. Under the preceding regularity assumptions this estimator is unbiased, and a cost-minimization update takes the form $\theta_{n+1} = \theta_n - \eta_n \widehat{\nabla_\theta J}(\theta_n)$ for a suitable step size $\eta_n > 0$. It is model-free in the relevant sense: it uses sampled costs and the policy score, but neither the transition kernels nor its derivatives.

The estimator may nevertheless have high variance. If $b_t(X_t)$ is any integrable baseline independent of $A_t$ conditionally on $X_t$, then \eqref{eq:background-score_zero_expectation} implies
\begin{align}
    \label{eq:background-reinforce_baseline}
    \nabla_\theta J(\theta) = \E^{\pi^\theta} \left[\sum_{t=0}^{T-1} \nabla_\theta\log p_t^\theta(\alpha_t\mid X_t) (C_t - b_t(X_t)) \right].
\end{align}
Subtracting the baseline therefore preserves unbiasedness while potentially reducing variance. A common choice is the conditional expected cost-to-go $b_t(x)=\E^{\pi^\theta}[C_t\mid X_t=x]$, although it is not in general the variance-minimizing baseline \cite{greensmith2004variance}. Estimating this quantity jointly with the policy leads to actor--critic methods. The present work instead builds on the score-function structure of REINFORCE.


\section{Mean-field control}
\label{sec:background-mfc}

Mean-field control (MFC) describes the cooperative control of a large population of interacting agents. In a system of $N$ agents, the interaction typically occurs through the empirical distribution $\mu_t^N = \frac1N\sum_{i=1}^N \delta_{X_t^{i,N}}$. The dynamics and cost of one agent may depend on its own state and action, but also on $\mu_t^N$. Directly treating the joint state $(X_t^{1,N},\ldots,X_t^{N,N})$ quickly becomes intractable as $N$ grows. The mean-field limit replaces the empirical distribution by the law $\mu_t=\Lc(X_t)$ of a representative agent. Under suitable assumptions, this yields a limiting control problem expressed in terms of the representative state and the population law, rather than the joint state of all $N$ agents \cite{motte2022mean}.

Mean-field control should be distinguished from mean-field games. In a mean-field game, each agent optimizes its own objective while treating the population flow as fixed, and an equilibrium condition subsequently requires the flow generated by the optimal response to coincide with the prescribed flow. In mean-field control, a central planner chooses a policy for the whole population and internalizes its effect on the population flow. The problems studied in this report are of the latter, cooperative type.

\subsection{Representative-agent formulation}

In discrete time, a randomized feedback policy may depend on both the individual state and the current population law:
\begin{align}
    \label{eq:background-mfc_dynamics}
    \begin{cases}
        X_0 &\sim \mu_0,\\
        \alpha_t &\sim \pi_t(\cdot\mid X_t,\mu_t),\\
        X_{t+1} &\sim P_t(\cdot\mid X_t, \mu_t,\alpha_t), \qquad t=0,\ldots,T-1,\\
        \mu_t &= \Lc(X_t).
    \end{cases}
\end{align}
The corresponding objective is
\begin{align}
    \label{eq:background-mfc_objective}
    J(\pi) = \E_{\mu_0,\pi}\left[ \sum_{t=0}^{T-1} r_t(X_t,\mu_t,\alpha_t) + g(X_T,\mu_T) \right],\qquad \inf_\pi J(\pi).
\end{align}
Unlike an exogenous state variable, the population law is not sampled independently of the policy. Once $\mu_t$ and $\pi_t$ are fixed, its next value is determined by
\begin{align}
    \label{eq:background-mfc_flow}
    \mu_{t+1}(B) = \int_{\Xc\times\Ac} P_t(B\mid x,a,\mu_t) \pi_t(\d a\mid x,\mu_t)\mu_t(\d x)
\end{align}
for every measurable set $B\subseteq\Xc$. Thus, changing the policy changes the whole sequence $(\mu_t)_{t=0}^T$, and this sequence in turn changes the policy, dynamics and costs evaluated at later times. This feedback is the defining feature of the problem.

\subsection{Lifted control problem}

At the population level, the state of the control problem is $\mu_t$ rather than $X_t$ alone. A control action at this level is a rule assigning an action distribution to each individual state. This gives a lifted MDP on $\Pc(\Xc)$, for which dynamic programming can be formulated when the model is known \cite{motte2022mean}. The lifted viewpoint also makes clear why the population law cannot simply be omitted from the state: two populations may contain an agent in the same state $x$ while inducing different transitions, costs and optimal decisions.

The difficulty is that the lifted state space is large. If $\Xc$ contains $d$ states, $\mu_t$ belongs to a $(d-1)$-dimensional simplex. For a continuous state space, $\Pc(\Xc)$ is infinite-dimensional.

\subsection{Representing the population law}

For a finite state space, the population law is represented exactly by a vector of state probabilities. There is no comparable canonical vector representation when the state space is continuous. In computations, one instead describes a measure $m$ through functions
\begin{align}
    \label{eq:background-measure_representation}
    \Psi(m) = \left(F_1(m),\ldots,F_d(m)\right).
\end{align}
such as moments, Fourier coefficients, interaction fields or kernel features. For example, one may take $F_j(m)=\langle\varphi_j,m\rangle=\int_\Xc\varphi_j(x)m(\d x)$. A sufficiently rich separating family of test functions characterizes the measure. A finite collection is exact when the problem has a corresponding finite-dimensional structure, and otherwise provides a task-dependent approximation.

A related approach is based on cylindrical representations $\Phi_J(m) = \left(\langle\varphi_1,m\rangle,\ldots,\langle\varphi_J,m\rangle \right)$,
which reduce an operator on probability measures to a function of finitely many measure functionals. Mekkaoui, Pham and Warin establish a universal approximation result of this form for operators on labelled conditional distributions and polynomial moment maps in their numerical implementation \cite{LO}.

This representation issue also enters the policy-gradient construction. After the population law is randomized, the mean-field argument $M_t^{\lambda,\theta}$ is itself a random measure. The finite-dimensional object whose density is studied is therefore
\begin{align}
    \label{eq:background-random_measure_representation}
    Z_t^{\lambda,\theta} = \Psi\left(M_t^{\lambda,\theta}\right),\qquad \Lc\left(Z_t^{\lambda,\theta}\right)\in\Pc(\R^d).
\end{align}
Thus the problem is not only to represent the population law, but to determine the joint law of functionals that characterize the relevant part of it. This law may admit an ordinary density and score even though there is no canonical Lebesgue measure on $\Pc(\Xc)$. The choice, or learning, of $\Psi$ is therefore part of the method: it must retain the distributional information used by the dynamics, costs and policy while remaining estimable from population samples.

\subsection{Model-free reinforcement learning for MFC}

MFC is closely related to cooperative multi-agent reinforcement learning (MARL). In a finite-agent formulation, learning a joint policy becomes difficult as the joint state and action spaces grow with population. Mean-field methods replace the individual interactions by aggregate population quantities, as in mean-field MARL \cite{yang2018mean}. For cooperative systems, the connection is made more precise in \cite{mondal2022approximation}, which studies how finite-agent MARL problems can be approximated by MFC problems. Mean-field MARL and MFC are not interchangeable, however the latter is the limiting control problem of a cooperative population under a common objective, which is the setting considered here.

When the transition and cost functions are unknown, the policy or value function must be learned directly from simulated or observed trajectories. Most model-free methods for MFC are value-based. A two-timescale mean-field $Q$-learning algorithm in discrete state and action spaces is proposed in \cite{angiuli2022unified}. In \cite{carmona2023model}, MFC is formulated as a lifted mean-field MDP and tabular and deep $Q$-learning methods are developed on this space. These methods avoid using the transition model explicitly, but the value function still takes a population distribution as part of its input.

Policy-based methods optimize a parametrized policy directly and are useful when the policy has a compact representation or the action space is continuous. Related actor--critic methods have been developed for continuous-time MFC with randomized policies \cite{frikha2025actor}, and more recently for extended MFC with deterministic policies \cite{cheng2026actor}. These approaches use continuous-time gradient representations and therefore differ from the discrete-time likelihood-ratio setting of this report. In discrete time, standard REINFORCE is not directly applicable: changing the policy parameters also changes the population flow $\mu_t^\theta$, hence the costs and transition kernels evaluated along that flow. The usual policy-score term does not account for this dependence.

The closest work is \cite{meunier2026model}, which introduces the first fully model-free policy-gradient method for discrete-time MFC with a finite state space. A strictly positive population vector is represented by its logits, Gaussian noise is added in logit space, and the perturbed vector is mapped back to the simplex by a softmax. A change of variables then replaces the unknown measure derivatives by an additional score term. The resulting MF-REINFORCE algorithm estimates the sensitivity of the population logits through an auxiliary simulation procedure. The authors prove convergence of the perturbed gradient as the perturbation vanishes and give explicit bias and mean-square error bounds for the estimator.

This construction is the starting point of the present work, but its logit representation is tied to finite state spaces: it requires one coordinate per state and has no direct counterpart on $\Pc(\Xc)$ when $\Xc$ is continuous. The method developed here separates the perturbation from the representation of the law. In the discrete case, the perturbation is applied directly on the simplex through a convex mixture, leading to an explicit simplex-density score and a recursive sensitivity estimator that reuses trajectories across the horizon. In the continuous case, the law is perturbed by probability-preserving push-forward maps, while the score is studied through the distribution features in \eqref{eq:background-random_measure_representation}. The extension is therefore not only a change of state space: it requires selecting or learning a representation of the population law and characterizing the law of that representation under the perturbation.
